\chapter{Watching the Conversation}

\section*{12.1}
One state corresponds to one residual finite session type, plus a terminal accepted state.

\section*{12.2}
The local trace must be complete and fully observed through the monitor's event abstraction.

\section*{12.3}
Acceptance/rejection classifies observed behavior; recovery requires actions and fault semantics not supplied by the automaton theorem.

\section*{12.4}
States expect `send Nat`, then `recv Bool`, then `close`, then accept.

\section*{12.5}
For a session beginning `send Nat`, a first event `recv Nat` is rejected at the initial state.

\section*{12.6}
`send Nat` alone is valid so far but incomplete because the residual session still requires `recv Bool; close`.

\section*{12.7}
Static typing reasons before execution from source/derivation structure. Monitoring reasons during execution from the events it can observe.

\section*{12.8}
A crash before instrumentation emits an event, or a semantic payload error that preserves type/order, may be invisible.

\section*{12.9}
The generated send-state has exactly one transition class, `send A`, to the monitor for the continuation. Acceptance after that transition reduces to the induction hypothesis.

\section*{12.10}
An internal-choice monitor has outgoing transitions only for the declared labels, each to the monitor generated from that continuation; apply the induction hypothesis on the chosen branch.

\section*{12.11}
Store observed event, expected class/set, residual session state, result, observation source, and timestamp separately.

\section*{12.12}
Distributed correlation needs an assumption such as trustworthy causal identifiers or synchronized clocks; without it, event order across peers may be ambiguous.
